\documentclass[11pt]{article}
\usepackage[utf8]{inputenc}
\usepackage{amsmath,amssymb}
\usepackage[margin=1in]{geometry}
\usepackage{hyperref}
\emergencystretch=3em

\title{The constant-Gaussian fluctuation model: the finiteness threshold}
\author{Claude, with Daniel Murfet}
\date{2026-09-07}

\begin{document}
\maketitle
\sloppy

\begin{abstract}
Second half of the rem:pop\_vs\_emp dichotomy. For $X\sim N(0,v)$ the expected Gaussian moment
$\mathbb E_+[J_p(X)]=\int_0^\infty s^{p-1}e^{-(\beta-\beta^2v/2)s^2}ds$ (unit~161) is
$J_p(0)(1-\beta v/2)^{-p/2}<\infty$ when $\beta v<2$ (\texttt{lintegral\_gaussMomentJ\_eq\_of\_lt}),
strictly larger than $J_p(0)$ for $v>0$ (\texttt{lt\_lintegral\_gaussMomentJ}), and $+\infty$ when
$\beta v\ge2$ (\texttt{lintegral\_gaussMomentJ\_eq\_top\_of\_ge}); hence
$\mathbb E_+[J_p(X)]<\infty\iff\beta v<2$ (\texttt{lintegral\_gaussMomentJ\_lt\_top\_iff}). The expected
leading coefficient can be infinite although every sample coefficient is finite. Zero
\texttt{sorry}/\texttt{axiom}.
\end{abstract}

\section{The things formalised}
{\small
\begin{verbatim}
theorem integral_rpow_mul_exp_neg_mul_sq (c p) (hc : 0 < c) (hp : 0 < p) :
    ∫ s in Ioi 0, s^(p-1) * exp (-c * s^2) = c^(-p/2) * (1/2) * Gamma (p/2)
theorem gaussMomentJ_zero (β p) (hβ) (hp) : gaussMomentJ β p 0 = β^(-p/2) * (1/2) * Gamma (p/2)
theorem gaussTail_eq_of_lt (β p v) (hβ) (hp) (hv : 0 ≤ v) (h : β * v < 2) :
    ∫⁻ s in Ioi 0, ofReal (s^(p-1) * exp (-β*s^2) * exp (v*(β*s)^2/2))
      = ofReal (gaussMomentJ β p 0 * (1 - β*v/2)^(-p/2))
theorem lintegral_rpow_Ioi_one_eq_top (p) (hp : 0 < p) : ∫⁻ s in Ioi 1, ofReal (s^(p-1)) = ⊤
theorem gaussTail_eq_top_of_ge (β p v) (hβ) (hp) (h : 2 ≤ β * v) :
    ∫⁻ s in Ioi 0, ofReal (…) = ⊤
theorem lintegral_gaussMomentJ_eq_of_lt / _eq_top_of_ge (X) (hX) (v : NNReal)
    (hXg : μ.map X = gaussianReal 0 v) …
theorem lintegral_gaussMomentJ_lt_top_iff ... :
    (∫⁻ ω, ofReal (gaussMomentJ β p (X ω)) ∂μ) < ⊤ ↔ β * v < 2
theorem lt_lintegral_gaussMomentJ ... (hv : 0 < v) (h : β * v < 2) :
    ofReal (gaussMomentJ β p 0) < ∫⁻ ω, ofReal (gaussMomentJ β p (X ω)) ∂μ
\end{verbatim}
}

\section{Mathematics}
Combining the exponentials, the integrand is $s^{p-1}e^{-cs^2}$ with $c=\beta(1-\beta v/2)$. For
$c>0$ the Gamma integral gives $c^{-p/2}\Gamma(p/2)/2$, and $J_p(0)=\beta^{-p/2}\Gamma(p/2)/2$, so the
ratio is $(1-\beta v/2)^{-p/2}$, which exceeds $1$ when $0<1-\beta v/2<1$. For $c\le0$ the integrand
dominates $s^{p-1}$ on $(1,\infty)$, whose integral diverges for $p>0$, so the Lebesgue integral is
$+\infty$; the extended-real formulation makes this a legitimate value rather than an integrability
failure.

\section{Paper-facing meaning}
rem:pop\_vs\_emp says the expected empirical coefficient is not the population coefficient. The
constant-Gaussian model sharpens this in two ways recorded in the mirror's remark: a Jensen gap
($\mathbb E[J_p(X)]>J_p(0)$) and a temperature threshold ($\beta v<2$) beyond which the expected
leading coefficient is infinite. As Astra~\#12 stressed, this concerns the expected asymptotic
coefficient; the fixed-$n$ chart integral on a compact chart has finite expectation in all regimes.
The model is a Gaussian constant fluctuation, not the empirical process itself (whose finite-$n$
exponential moments need not exist).

\section{Lean notes}
\begin{itemize}
\item \texttt{integral\_rpow\_mul\_exp\_neg\_mul\_rpow} with exponent $2$ and \texttt{Real.rpow\_two}
converts to the \texttt{Monoid} power $s^2$ used in our integrands.
\item Non-integrability of $s^{p-1}$ on $(1,\infty)$ is \texttt{integrableOn\_Ioi\_rpow\_iff one\_pos}
(iff $p-1<-1$); converting to $\int^-=\top$ is \texttt{hasFiniteIntegral\_iff\_ofReal} with
\texttt{not\_lt} and \texttt{top\_le\_iff}, using that \texttt{Integrable} is the pair
(\texttt{AEStronglyMeasurable}, \texttt{HasFiniteIntegral}).
\item Pointwise comparison must be restricted to $s>1$ (\texttt{lintegral\_mono\_ae} with
\texttt{ae\_restrict\_of\_forall\_mem}); a global \texttt{lintegral\_mono} fails because \texttt{rpow}
with negative base is not controlled.
\item \texttt{Real.one\_lt\_rpow\_of\_pos\_of\_lt\_one\_of\_neg} for the Jensen gap;
\texttt{ENNReal.ofReal\_lt\_ofReal\_iff} needs positivity of the right argument.
\end{itemize}

\end{document}
